\subsection{Strict positivity from the tautological metric}
\label{subsec:tautological-homogeneous-positivity}

Let $F\to X$ be a holomorphic vector bundle of positive rank over a compact complex manifold. In this subsection, $\mathbb P(F)$ parametrizes complex lines in $F$, and
\[
 L=\mathcal O_{\mathbb P(F)}(1)
   =\mathcal O_{\mathbb P(F)}(-1)^*.
\]
The complement $F^\times$ of the zero section is holomorphically identified with the complement of zero in the tautological line bundle by $v\mapsto([v],v)$. We use the curvature convention $R_h=-\ddc\log h$.

\begin{lemma}\label{lem:tautological-homogeneous-levi}
A smooth Hermitian metric $h$ on $L$ determines a positive smooth complex homogeneous function $H$ of degree two on $F^\times$, by taking the squared norm in its dual metric on the tautological line. The function $H$ is strictly plurisubharmonic if and only if $R_h$ is strictly positive on $\mathbb P(F)$.
\end{lemma}
\begin{proof}
Take local holomorphic coordinates $y$ on $\mathbb P(F)$ and a holomorphic frame $e$ of the tautological line. Write the metric of its dual frame as $h=\mathrm e^{-\varphi}$. The dual metric gives
\[
 H(y,\lambda e)=|\lambda|^2\mathrm e^{\varphi(y)},\qquad
 R_h=\ddc\varphi.
\]
For the complex Hessian $\mathcal L$, without the $2\pi$ curvature factor, differentiation at $\lambda\ne0$ gives
\[
 \mathcal L H(u,t)
 =H\left(\mathcal L\varphi(u)+
          \left|\frac{t}{\lambda}+\partial\varphi(u)\right|^2\right),
 \qquad (u,t)\in T_y^{1,0}\mathbb P(F)\oplus\mathbb C.
\]
If $\mathcal L\varphi$ is positive definite, the right-hand side is positive for every nonzero $(u,t)$: the first term detects $u\ne0$, and the second detects $t\ne0$ when $u=0$. Conversely, set $t=-\lambda\partial\varphi(u)$ for $u\ne0$. Strict positivity of $\mathcal L H$ then forces $\mathcal L\varphi(u)>0$. This is exactly positivity of $R_h$. The squared tautological norm is independent of the chosen frame and is complex homogeneous, so the local functions define the asserted global $H$.
\end{proof}

A line bundle is very ample when its complete section map is a holomorphic embedding; we use ampleness in the form that some positive tensor power is very ample. In the line convention above, $F^*$ is ample when $L$ is ample.

\begin{proposition}\label{prop:ample-dual-strict-homogeneous-function}
If $F^*$ is ample, there is a positive complex homogeneous function $H$ of degree two, smooth and strictly plurisubharmonic on $F^\times$ and continuous on $F$ after extension by zero.
\end{proposition}
\begin{proof}
Choose $m\geq1$ such that $L^m$ is very ample. Its complete section map gives an embedding $\iota:\mathbb P(F)\to\mathbb P^N$ with $L^m=\iota^*\mathcal O_{\mathbb P^N}(1)$. Pull back the Fubini--Study metric to a metric $k$ on $L^m$, and take its positive $m$th root $h$ on $L$. In a local frame this means $h=k^{1/m}$; the transformation law for $L^m$ makes these roots a genuine metric on $L$. Its curvature is
\[
 R_h=\frac1m\iota^*\omega_{\rm FS}>0,
\]
because the differential of the embedding is injective. Apply \cref{lem:tautological-homogeneous-levi}. The resulting $H$ is bounded on the unit sphere bundle of a fixed Hermitian metric, since the base is compact. Homogeneity then gives $0<H(v)\leq C|v|^2$ off zero, which proves continuity of the extension by zero.
\end{proof}

\begin{corollary}\label{cor:ample-dual-continuous-convex-obstruction}
Let $F$ carry a holomorphic structure in the neighbourhood of \cref{thm:convex-norm-deformation-obstruction}. If $F^*$ is ample, there is no positive continuous complex homogeneous function of degree two on $F$ that is plurisubharmonic off zero and locally strongly convex in the fibres.
\end{corollary}
\begin{proof}
By \cref{prop:ample-dual-strict-homogeneous-function}, dual ampleness supplies the function $H$ required in \cref{prop:weak-homogeneous-convex-replacement}. A function $Q$ as in the statement would therefore give a smooth complex norm on $F$ with strictly plurisubharmonic square. This contradicts \cref{thm:convex-norm-deformation-obstruction}.
\end{proof}

The tautological construction permits the homogeneous function to depend on the fibre line, and supplies the strict total-space positivity used in regularization. Local strong fibre convexity remains an independent hypothesis in the continuous obstruction.
